% =====================================================================
%  §1 Introduction  （原文 p1–p4 开头）
% =====================================================================
\section{引言}

\subsection{同调镜像对称与退化}

数学上，镜像对称可以有多种解释。本文要在两种进路之间架一座桥：同调镜像对称
（homological mirror symmetry，见 \cite{Ko}），以及环面纤维化之间的对偶（这是
Strominger--Yau--Zaslow 猜想的一个版本，见 \cite{SYZ}）。

镜像对称是 Calabi--Yau 流形之间的对偶，所谓 Calabi--Yau 流形，即带有 Ricci 曲率消失的
凯勒度量的复流形。事实上，它与其说是单个流形之间的对偶，不如说是处于某种「退化」族中的
流形之间的对偶（「大复结构极限」与「大辛结构极限」）。本文提出这一退化的一个微分几何模型。
特别地，我们猜想：在极限下，两个对偶流形 $X$ 与 $X^{\vee}$ 都成为以环面为纤维的纤维丛，
而且底空间是同一个 $\bar Y$（见 \S3）。度量空间 $\bar Y$ 是某个（实）黎曼流形 $Y$ 的紧化，
带有一些温和的奇点，而 $Y$ 的维数是 $X$ 与 $X^{\vee}$ 维数的一半。此外，流形 $Y$ 还带有
一个丰富的几何结构，其中包含某种「组合」数据（即所谓的整仿射结构，integral affine structure）。这幅图景部分地受到
M.~Gromov 等人发展的坍缩黎曼流形理论的启发（例如见 \cite{CG}）。我们这些几何猜想的另一个
来源是 \cite{SYZ} 版本的镜像对偶。我们在 \S2 中用略有不同的语言把它重新表述一遍，那一节
讨论共形场论的模空间及其自然的紧化。在最近的一篇预印本 \cite{GW} 中，类似的微分几何猜想被
提出，并在退化的 K3 曲面的情形得到了验证。

由 \cite{Ko} 提出的同调镜像猜想，是关于两个 $\Ainf$-范畴等价的一个陈述：
Calabi--Yau 流形 $X$ 上（凝聚层的导出）范畴，与对偶的 Calabi--Yau 流形 $X^{\vee}$ 的
Fukaya 范畴。前者用全纯（或代数）的语言定义，后者则用辛几何的语言定义。

我们把坍缩黎曼流形理论的想法应用到同调镜像猜想上。我们把退化 Calabi--Yau 流形的这个模型
称为\emph{微分几何}模型。它给出了 Fukaya 范畴退化的清晰图景。辛流形 $(X^{\vee},\omega^{\vee})$
（其中 $[\omega^{\vee}]\in H^2(X^{\vee},\Z)$）的 Fukaya 范畴 $F(X^{\vee},\omega^{\vee})$ 及其退化，
都定义为 Laurent 形式幂级数域 $\C((q))$ 上的 $\Ainf$-范畴。参数 $q$ 是这样进入的：当我们写
高阶复合（Massey 积）时，其系数形如 $q^{\int_{\beta}\omega^{\vee}}$，$\beta\in H_2(X^{\vee},\Z)$。
可以令 $q=\exp(-1/\eps)$，$\eps\to 0$，其中参数 $\eps$ 对应于辛形式的重新标度：
$\omega^{\vee}\mapsto \omega^{\vee}/\eps$。如果 $[\omega^{\vee}]$ 不属于 $H^2(X^{\vee},\Z)$，
就可以在域
\[
  \C_{\eps}:=\Bigl\{\sum_{i\ge 0}a_i\,e^{-\lam_i/\eps}\ \Bigm|\ a_i\in\C,\ \lam_i\in\R,\ \lam_i\to+\infty\Bigr\}
\]
上工作。

在环面纤维化的情形，极限 Fukaya 范畴的一个满子范畴可以用底空间上的 Morse 理论来描述。
给出 $\Ainf$-结构的高阶乘法，可以写成对一族平面树求和。在余切丛的情形（而不是环面纤维化的
情形），这一描述早先已由 Fukaya 与 Oh 提出（见 \cite{FuO}）。

在镜像对称的全纯一侧，对偶族 $X_q$ 的退化用非阿基米德的语言描述：我们有一个 Calabi--Yau
流形 $\Xmer$，定义在 $q=0$ 处亚纯函数芽之域 $\Cmer$ 上；改变标量，便得到 Laurent 级数
局部域 $\C((q))$ 上的 Calabi--Yau 流形 $\Xform$。我们把退化
的这一图景称为\emph{解析}图景。对于（任意局部域上、关于离散赋值完备的）代数 Calabi--Yau
流形的一类，存在用带整仿射结构的实 $C^{\infty}$-流形给出的描述。我们期待退化
的微分几何图景与非阿基米德解析图景是等价的。这一等价反映了看待 Calabi--Yau 流形的两种不同
方式：微分几何的方式（通过 Ricci 曲率消失的凯勒度量）与代数几何的方式（通过典范类消失的
光滑射影簇）。

同调镜像猜想说，Fukaya 范畴 $F(X^{\vee},\omega^{\vee})$（作为 $\C((q))$ 上的 $\Ainf$-范畴）
等价于凝聚层的导出范畴 $D^b(\Xform)$。我们期待它蕴涵关于 Calabi--Yau 流形上有理曲线条数
（亏格零 Gromov--Witten 不变量）的若干众所周知的数值预言。

利用我们关于 Calabi--Yau 流形坍缩的猜想，我们给出证明同调镜像猜想的一条一般进路。我们在
环面纤维化没有奇点的情形应用它，这出现在阿贝尔簇的情形。一般情形下，应当研究环面纤维化的
底空间之奇点所带来的影响。沿着同样的思路，我们在平坦环面与阿贝尔簇这一最简单的情形中，
讨论带整仿射结构的黎曼流形与非阿基米德域上的簇之间的关系。一般情形将另文讨论。

从上面的讨论应当已经清楚，非阿基米德分析在同调镜像猜想的表述与证明中起着重要作用。
退化的解析图景似乎与 Berkovich 版本的刚性解析空间理论有关（见 \cite{Be}）。特别地，
环面纤维化的底空间与局部域上代数 Calabi--Yau 流形的 Berkovich 谱中某个典范地定义的子集
（见 \S3.3）之间，有着惊人的相似。这个课题无疑值得进一步研究。

\begin{yizhu}[同调镜像猜想的出处]
同调镜像猜想是 Kontsevich 在 1994 年国际数学家大会的报告 \cite{Ko} 中提出的。其原始表述是
\[
  D^b(\Coh(X))\ \simeq\ D^{\pi}\Fuk(X^{\vee}),
\]
即全纯一侧取凝聚层的导出范畴，辛一侧取 Fukaya 范畴（可能须先作三角化和 Karoubi 完备化）。
本文把两侧都换成 $\Ainf$-范畴、并取 $\C((q))$ 为其系数域，正是为了把「退化」这一极限过程
纳入范畴的代数结构之中：$q$ 记录辛形式的标度，$q\to 0$ 就是大辛结构极限。
\end{yizhu}

\begin{yizhu}[整仿射结构]
所谓流形 $Y$ 上的整仿射结构，是指在 $Y$ 的每点的切空间上给定一个仿射结构，使转移映射落在
$\mathrm{SL}(n,\Z)\ltimes \R^n$ 中；等价地说，$Y$ 有一个图册，其转移映射是整系数的仿射变换。
它是「组合」数据的来源：局部常值的整切标架、以及底流形上的 $\Z$-仿射几何。本文 \S3.2 会
说明它如何由 Monge--Ampère 流形的对偶纤维化产生。
\end{yizhu}

\subsection{本文内容}

在 \S2 中我们讨论来自共形场论的动机。在 \S3 中我们表述关于大复结构极限的解析图景与微分几何
图景的猜想。在 \S4 中我们描述一套 $\Ainf$-预范畴的一般框架，它适应于 Fukaya 范畴定义中的
横截性问题。\S5 讨论 Fukaya 范畴及其退化。读者会注意到在 Laurent 幂级数域上工作的一个好处：
可以考虑拉格朗日子流形上的\emph{所有}局部系（在通常的进路中则要求和乐的酉性）。\S6 讨论
由 Fukaya 引入的（随后由 Fukaya 与 Oh 在 \cite{FuO} 中研究的）光滑函数的 $\Ainf$-范畴。
我们证明这个 $\Ainf$-范畴有一个非常简单的 de Rham 模型。本文的这一部分可以独立于其余部分
阅读。另一方面，其中的技巧与证明的一般方案，稍后将在 \S8 中用到，那一节专门讨论同调镜像
猜想。一个重要的技术工具是微分分次代数的子复形上的显式 $\Ainf$-结构（见 \cite{GS}、
\cite{Me}）。我们把 \cite{Me} 中的公式用对一族平面树求和的语言重新表述。我们关于 Morse
与 de Rham 两个 $\Ainf$-范畴等价的证明，采用了 \cite{HL} 中处理 Morse 理论的方式。
\S7 讨论同调镜像猜想的解析方面。我们给 \S3 中讨论的那一类环面纤维化指定一个刚性解析空间。
我们还用非阿基米德几何的语言，为环面纤维化构造一个镜像对称函子。使用非阿基米德分析，
使我们得以避开 Fukaya 范畴定义中级数收敛性的问题。在 \S8 中我们构造一个 $\Ainf$-预范畴，
它等价于 $\C_{\eps}$ 上 Calabi--Yau 流形的凝聚层导出范畴的 $\Ainf$-版本的一个满
$\Ainf$-子范畴。我们证明这个范畴等价于镜像对偶的环面纤维化的 Fukaya 范畴的某个
$\Ainf$-子范畴。在附录（\S9）中，我们描述这些构造在复几何情形下的类似物。

\medskip
\noindent\textbf{致谢.}\quad Y.S.~的研究部分地得到 IHES、Clay 数学研究所的 fellowship 以及
CRDF 基金 UM-2091 的资助。他也感谢 IHES 的款待与优良的研究条件。

\begin{yizhu}[关于 $\C((q))$ 与 $q=\exp(-1/\eps)$]
Laurent 形式幂级数域 $\C((q))=\C[[q]][q^{-1}]$ 在这里扮演「形式化的大辛结构极限」的角色。
把 $q$ 取成 $\exp(-1/\eps)$ 之后，$q^{\int_\beta\omega^{\vee}}=e^{-\frac{1}{\eps}\int_\beta\omega^{\vee}}$，
正是路径积分中瞬子（全纯曲线）贡献的典型形式：$\eps\to0$ 即大辛结构极限。若 $[\omega^{\vee}]$
不是整系数的，则 $\int_\beta\omega^{\vee}$ 不是整数，$q$ 的指数取值于实数，此时需把系数域
扩大为 $\C_{\eps}$，这正是文中紧接着引进 $\C_{\eps}$ 的原因。
\end{yizhu}

\begin{yizhu}[原文的一处笔误]
原文此处写作「If $[\omega^{\vee}]$ does not belong to $H^2(X,\Z)$」，其中 $X$ 应为
$X^{\vee}$：讨论的对象是辛侧流形 $X^{\vee}$ 的辛类，$X$ 是复侧流形。译文已按
$X^{\vee}$ 译出。
\end{yizhu}
